% !TeX root = ../main.tex
论文基于MOOSE有限元框架，开展了快堆MOX燃料微观结构演化、多物理场耦合建模、燃料棒性能验证及POD–LSTM快速预测方法研究，论文选题具有重要的科学意义和工程应用价值。

论文取得的主要创新性成果包括：

\begin{enumerate}[wide]
    \item 建立了微观结构演化与多组分迁移耦合的快堆MOX燃料分析有限元框架，揭示了气孔增强输运的空间局部化特征及其对组分再分布的作用机制；
    \item 建立了材料变形、裂变气体行为、气腔压力及燃料–包壳间隙传热之间的耦合计算方法，通过数值基准和公开辐照实验数据验证了相关模型的准确度；
    \item 提出了基于POD–LSTM的氧化物燃料元件热–力场快速预测方法，实现了燃料元件二维完整热–力场的快速重构，显著提高了计算效率。
\end{enumerate}

论文写作规范，态度科学严谨，数据详实，工作量饱满。论文工作表明作者在本学科具有坚实宽广的基础理论和系统深入的专门知识，具备独立从事科学研究工作的能力。答辩过程中，条理清晰，叙述清楚，回答问题正确。

答辩委员会经无记名投票表决，5票一致同意通过博士论文答辩，并建议授予王召浩工学博士学位。
